%-*-latex-*-
%-*-latex-*-
\newcommand\COURSE{ciss350}
\newcommand\ASSESSMENT{q10801}
\newcommand\ASSESSMENTTYPE{Quiz}
\newcommand\POINTS{\textwhite{xxx/xxx}}

\input{myquizpreamble}
\input{yliow}
\input{\COURSE}
%-*-latex-*-

%-*-latex-*-

%-*-latex-*-
\renewcommand\TITLE{\ASSESSMENTTYPE \ \ASSESSMENT}

\renewcommand\EMAIL{}
\renewcommand\AUTHOR{eagraham1@cougars.ccis.edu}

\textwidth=6in
\begin{document}
\topmatterthree

For the following questions,
the following is how you describe a hashtable.
Suppose you have a collection of information on
people t-shirt size:
\begin{console}[fontsize=\footnotesize]
John, L
Mary, L
Sue, S
Tom, XL
\end{console}
and you are using a table of size 5, and
the mod 5 values of the hash of John, Mary, Sue, Tom are
1, 0, 2, 4, you write
you write
\begin{console}[fontsize=\footnotesize]
[0, Not-Available, Mary, L]
[1, Not-Available, John, L]
[2, Not-Available, Sue, S]
[3, Available, None, None]
[4, Not-Available, Tom, XL]
\end{console}
The first column is the index value, the second is the hashtable row flag
which must be Available, Not-Available, or Deleted,
and third column is the key (the person's name),
and the fourth column is the t-shirt size.
The flag for each row must be set to one of the above three values.
The key column and value column (third and fourth column) are initialized
to None.
If there's collision, linear probing is used.

TURN PAGE ...

%------------------------------------------------------------------------------
\newpage
\nextq
You work for NSA.
A message in characters a, b, c, d, e, f, g, h has been encrypted into 4-bit
codes, i.e.,
each 4-bit sequence corresponding to one of the characters
a, b, c, d, e, f, g, h.
You do know the frequencies of a, b, c, d, e, f, g, h for the original message.
For instance c is the most commonly occurring character.
So you want to build a frequency table.
Here's an example:
\begin{console}[fontsize=\footnotesize]
0000, 5
0010, 6
0100, 3
0110, 8
1000, 10
1011, 12
0111, 3
\end{console}
In this example 0000 occurs 5 times in the message, 0010 occurs 6 times in the
message, etc.
The above are the only possible 4-bit patterns in the message.
Create a hashtable to build the frequency table of the following
message:
\[
0010, 0000, 0100, 0110, 0111, 0010, 0000, 1000, 0100, 1011, 0111, 0100, 1000
\]
Assume the hash function on those bit patterns give you the following:
\begin{console}[fontsize=\footnotesize]
key   hash(key)
0000  12
0010  3
0100  7
0110  15
1000  22
1011  12
0111  42
\end{console}
and assume that you are using a hashtable of size 20.
Write down the resulting hashtable.
\\
\ANSWER
\begin{answercode}
Key   Hash  Index  Frequency
0000  12    12     2
0010  3     3      2
0100  7     7      3
0110  15    15     1
1000  22    2      2
1011  12    12     1
0111  42    2      2
\end{answercode}

%------------------------------------------------------------------------------
\nextq
In the above, how many collisions did you encounter?
\\
\ANSWER
\begin{answercode}
2 collsions
\end{answercode}

%------------------------------------------------------------------------------
\nextq
Which key search will result in the longest chain of probes?
If there's more than one, write down all the bit patterns separated by
comma.
\\
\ANSWER
\begin{answercode}
1000, 0111, 1011, 0000
\end{answercode}

%------------------------------------------------------------------------------
\nextq
You have just been informed that the message was corruped during
transmission.
All the 0000 should have been 0110.
Write down the resulting new hashtable.
\\
\ANSWER
\begin{answercode}
Key   Hash  Index  Frequency
0010  3     3      2
0100  7     7      3
0110  15    15     3
1000  22    2      2
1011  12    12     1
0111  42    2      2
\end{answercode}

%------------------------------------------------------------------------------
%-*-latex-*-
\newpage
\input{instructions}
\end{document}

